\section{Giriş}\label{sec:giris}

Gezgin satıcı problemi (traveling salesman problem, TSP), rotalama, devre yerleşimi ve üretim
planlamasındaki pek çok kararın ortak çekirdeğidir \citep{reinelt_1991_tsplib, applegate_2006_tsp}.
Pratikte büyük örnekler yerel aramayla çözülür. Bu ailenin en yaygın iskeleti iteratif yerel
aramadır (iterated local search, ILS): tur bir yerel optimuma indirilir, ardından bir pertürbasyonla
sarsılır ve yeniden onarılır \citep{martin_1992_largestep, lourenco_2010_ils}. TSP'de bu pertürbasyon
çoğunlukla double-bridge vuruşudur (kick). Vuruşun hangi tur kenarlarını kıracağı, aramanın sınırlı
zamanı nereye harcayacağını belirler. Bu nedenle küçük görünen bu seçim, kısa bütçeli bir çözücünün
kalitesini doğrudan etkiler \citep{johnson_1997_case, applegate_2003_clk}.

\begin{figure}[!htbp]
\centering
% MANSET (hero) -- OTOMATIK URETIM: evidence/scripts/build_v3.py (figures/scripts/manset.py)
% Veri: figures/data/manset_diff_gnn.dat, manset_diff_random.dat, manset_vitrin.csv (vitrin_sec.py, kategori A)
% Kanit: E789 (vs GNN), E790 (vs Random); vitrin: evidence/vitrin.json
% vitrin_sec kazanma orani (gurultu |d|>0.0156 pt): 70.4%
\begin{tikzpicture}
\begin{groupplot}[group style={group size=1 by 2, vertical sep=1.25cm},
  width=\textwidth, height=5.4cm, xmin=0, xmax=55, xtick=\empty,
  xlabel={Holdout2 instances, sorted by $\Delta$ (54)}, title style={font=\small, yshift=-1ex},
  tick label style={font=\scriptsize}, label style={font=\small},
  y tick label style={/pgf/number format/fixed, /pgf/number format/precision=1},
  ymajorgrids, grid style={black!10}]
\nextgroupplot[ymin=-1.4, ymax=0.8,
  ylabel={$\Delta$ gap, Excess $-$ GNN [pt]}, title={(a) Excess vs.\ GNN}]
\addplot[ybar, bar width=0.72, fill=iyiRenk, draw=none] table[x=rank, y=win, unbounded coords=jump] {figures/data/manset_diff_gnn.dat};
\addplot[ybar, bar width=0.72, fill=kotuRenk, draw=none] table[x=rank, y=loss, unbounded coords=jump] {figures/data/manset_diff_gnn.dat};
\draw[black!60] (axis cs:0,0) -- (axis cs:55,0);
\node[anchor=north west, font=\small, fill=white, inner sep=2pt, align=left] at (axis description cs:0.01,0.97) {\textcolor{iyiRenk}{\textbf{Excess lower gap: 41/54}}\quad \textcolor{kotuRenk}{\textit{GNN lower gap: 13/54}}};
\node[circle, fill=black, inner sep=1.1pt] (v0) at (axis cs:1,-1.18350) {};
\node[anchor=west, font=\scriptsize, fill=white, inner sep=1pt, draw=black!30] (l0) at (axis cs:19,-1.0751) {h2c02 (clustered): Excess 0.29\%, GNN 1.47\%};
\draw[black!70, thin, -{Latex[length=1.5mm]}] (l0.west) -- (v0);
\node[circle, fill=black, inner sep=1.1pt] (v1) at (axis cs:3,-0.81961) {};
\node[anchor=west, font=\scriptsize, fill=white, inner sep=1pt, draw=black!30] (l1) at (axis cs:19,-0.7861) {eg7146 (real): Excess 1.53\%, GNN 2.35\%};
\draw[black!70, thin, -{Latex[length=1.5mm]}] (l1.west) -- (v1);
\node[circle, fill=black, inner sep=1.1pt] (v2) at (axis cs:17,-0.12025) {};
\node[anchor=west, font=\scriptsize, fill=white, inner sep=1pt, draw=black!30] (l2) at (axis cs:19,-0.4970) {h2u12 (uniform): Excess 0.46\%, GNN 0.58\%};
\draw[black!70, thin, -{Latex[length=1.5mm]}] (l2.west) -- (v2);
\node[anchor=east, font=\scriptsize, fill=white, inner sep=1pt, draw=black!30, text=kotuRenk] (c0) at (axis cs:51.5,0.2803) {counter-example h2c13 (clustered): Excess 2.16\%, GNN 1.53\%};
\draw[black!70, thin, -{Latex[length=1.5mm]}] (c0.east) -- (axis cs:53.6,0.2803);
\nextgroupplot[ymin=-1.5, ymax=0.6,
  ylabel={$\Delta$ gap, Excess $-$ Random [pt]}, title={(b) Excess vs.\ Random}]
\addplot[ybar, bar width=0.72, fill=iyiRenk, draw=none] table[x=rank, y=win, unbounded coords=jump] {figures/data/manset_diff_random.dat};
\addplot[ybar, bar width=0.72, fill=kotuRenk, draw=none] table[x=rank, y=loss, unbounded coords=jump] {figures/data/manset_diff_random.dat};
\draw[black!60] (axis cs:0,0) -- (axis cs:55,0);
\node[anchor=north west, font=\small, fill=white, inner sep=2pt, align=left] at (axis description cs:0.01,0.97) {\textcolor{iyiRenk}{\textbf{Excess lower gap: 44/54}}\quad \textcolor{kotuRenk}{\textit{Random lower gap: 10/54}}};
\end{groupplot}
\end{tikzpicture}

\caption{54 holdout2 probleminde 10~sn gap farkı, sıralı: (a) fazla-uzunluk$-$GNN, (b) fazla-uzunluk$-$rastgele; mavi: vekil lehine. Vitrinler, GNN'nin kendi çeyrekler arası aralığında kaldığı problemlerden katman başına en büyük avantajla seçildi; ölçüm gürültüsünün altındaki farklar beraberlik sayıldığında vekil problemlerin \%70'inde GNN'den iyidir.} % V013
\label{fig:manset}
\end{figure}

Vuruşu rastgele değil de ``şüpheli'' kenarlara yöneltmek doğal bir fikirdir. Öğrenilmiş TSP
yöntemlerinin başarılı bir kolu da tam olarak bunu yapar. Güçlü bir yerel aramanın içine bir model
yerleştirir ve modelin kenar tahminleriyle aramayı yönlendirir: grafik sinir ağı (graph neural
network, GNN) ile güdümlü yerel arama \citep{hudson_2022_gnngls}, öğrenilmiş aday kümeleri
\citep{xin_2021_neurolkh} ve ısı haritası güdümlü arama \citep{qiu_2022_dimes, sun_2023_difusco}
bunun örnekleridir \citep{karimimamaghan_2022_mlmeta}. Bu çizgiyi izleyerek,
Lin--Kernighan--Helsgaun (LKH) çözücüsünün LKH-3 sürümüyle üretilen turlarla eğitilmiş bir GNN'nin turda kalmayacağını öngördüğü kenarları vuruşla kırmak
kolay bir adımdır \citep{helsgaun_2017_lkh3}.

Bu adımın değerini ölçmek ise sanıldığı kadar kolay değildir. Öğrenilmiş bir rehber çoğunlukla
rastgele seçime ya da kendi varyantlarına karşı sınanır; aynı geometriden türetilebilecek en basit
eğitimsiz vekile karşı sınanması seyrektir. Maksimum bağımsız küme probleminde GNN çıktısının
rastgele değerlerle değiştirilmesinin başarımı pek değiştirmediği gösterilmiştir
\citep{bother_2022_wrong}. Büyük ölçekli TSP'de de uzaklıktan türetilen basit bir ısı haritası,
öğrenilmiş olanlarla rekabet edebilmiştir \citep{xia_2024_rethinking}. İkinci bir sorun, tahmin
doğruluğunun arama faydasıyla eş tutulmasıdır. Bir modelin ``yanlış'' kenarları iyi ayırt etmesi,
onları kırmanın turu kısaltacağı anlamına gelmez. Modelin çalışma süresi de aynı bütçeden
ödenmelidir. Son olarak, yöntem ve ayar seçimi ile doğrulama aynı veride yapıldığında olumlu sonuç
kolayca abartılır \citep{hooker_1995_testing, nosek_2018_prereg}.

Bu çalışmada vuruş hedeflemeyi bu üç tuzaktan arındırılmış bir düzenekte inceledik. Önce saniyede
binlerce vuruş atabilen hızlı bir ILS test ortamı kurduk ve bütün model sürelerini bütçeye yükledik.
Ardından çalışmayı üç aşamada yürüttük. Aşama~1'de yavaş bir iskelette gözlediğimiz anlamlı yararın
ölçüm gürültüsü olduğunu gördük ve bu iddiayı geri çektik. Aşama~2'de, ön-kayıtlı bir protokolle 38
yeni çok büyük ölçekli tümleşik devre (very-large-scale integration, VLSI) probleminde GNN
hedeflemesini rastgele vuruşa ve eğitimsiz bir $k$-en yakın komşu (kNN) sırasına karşı sınadık. % E797
Aşama~3 doğrulayıcı aşamadır. Yöntem ve havuz ayarları yalnız geliştirme setlerinde seçildi. Protokol
SHA-256 özetiyle donduruldu ve özet koşudan önce OpenTimestamps'a gönderildi; Bitcoin blok
976258'e bağlanan bu kanıt, protokolün sonuçlar analiz edilmeden önce var olduğunu üçüncü
kişilerce doğrulanabilir kılar. Test, geliştirmede hiç kullanılmamış 54 problemde (holdout2)
koşuldu. % E584, E585, E586

Sınanan adaylardan biri, güdümlü yerel aramanın (guided local search, GLS) ``uzun kenar şüphelidir''
ilkesini \citep{voudouris_1999_gls} vuruş hedeflemeye taşıyan, eğitimsiz ve kenar başına $O(1)$
maliyetli bir fazla-uzunluk vekilidir. Kenar uzunluğuna dayanan basit bir sinyaldir: uzunluktan uç
noktaların yerel komşu ölçeğini çıkarır. Şekil~\ref{fig:manset} sonucu problem düzeyinde gösterir: her çubuk bir problemdir ve
sıfırın altındaki mavi çubuklar vekilin daha düşük gap verdiği problemlerdir. Vekil 54 problemin
41'inde GNN'den, 44'ünde rastgele vuruştan daha iyi sonuç verir; çubukların en uzunları kümeli ve
gerçek problemlerdedir. % E789, E790
GNN'ye karşı ön-kayıtlı iddia bir kötü-değil iddiasıdır; problemlerin çoğundaki üstünlük keşifsel
bir bulgudur. Aynı ölçütü kendi vekilimize de uyguladık: protokol sonrası keşifsel bir kolda yerel
ölçek düzeltmesi olmadan ham kenar uzunluğu da vekilden ortalamada ayrışmadı. Bu, yararın kenar
uzunluğu bilgisinin kendisinden geldiğini düşündürür (Bölüm~\ref{sec:bulgular-duyarlilik}). % E972

Bu çalışmanın başlıca katkıları şunlardır:
\begin{itemize}
  \item Kenar uzunluğuna dayalı eğitimsiz fazla-uzunluk vekilinin geliştirmede kullanılmamış 54 problemde rastgele vuruşa göre ortalama
  gap'i 0,172 puan düşürdüğünü ($p=3{,}7\times10^{-7}$) ve GNN ile kazanç hedefli çok katmanlı
  algılayıcıya (multilayer perceptron, MLP) 0,05 puanlık sınırla kötü olmadığını ön-kayıtlı
  testlerle gösteriyoruz. % E600, E601, E602, E589
  \item Öğrenilmiş skorlarda kestirim doğruluğunun arama faydasına kendiliğinden dönüşmediğini
  ölçüyoruz: Aşama~2'de GNN yanlış kenarları 58 problemin 57'sinde kNN'den isabetli ayırt ettiği
  hâlde aramada onun gerisinde kaldı; holdout2'de ikisi ayrışmadı, ucuz ve sürekli güncel vekil ise
  ikisinden de iyi sonuç verdi (keşifsel). % E890, E805, E605, E604, E607
  \item Koşu içinde çevrimiçi öğrenmenin aynı doğrusal modelin sabit hâline göre gap'i 0,291 puan
  artırdığını ($p=1{,}6\times10^{-10}$) ve mekanizmasını ölçüyoruz: model kendi vuruş ödülünü
  yükseltirken (12/12 problem) yerel aramanın onardığı nötr vuruşlara yönelir (\%47'den \%88'e). % E599, E941, E944
  \item Geliştirmede seçim, SHA-256 ile dondurma, dış zaman damgası ve 54 bağımsız problemle
  doğrulamadan oluşan bir ön-kayıt protokolünü ve ``öğrenilmiş skoru önce en basit eğitimsiz vekile
  karşı sına'' ilkesini sunuyor, ilkeyi kendi vekilimize de uyguluyoruz: daha basit ham kenar
  uzunluğu keşifsel olarak vekilden ayrışmadı. % E585, E586, E972
\end{itemize}

Uygulayıcı için sonuç sadedir: bu motorda ve 10~sn bütçede vuruşları hedeflemek kazandırır;
sınanan öğrenilmiş skorlar karşısında bunun için model eğitmek gerekmedi. % E588

Makalenin geri kalanı şöyle düzenlenmiştir.
Bölüm~\ref{sec:ilgili-calismalar} ilgili çalışmaları, Bölüm~\ref{sec:yontem} test ortamını ve
hedefleme sinyallerini, Bölüm~\ref{sec:deney-kurulum} üç aşamanın protokollerini verir.
Bölüm~\ref{sec:bulgular} doğrulayıcı sonuçları ve mekanizmaları, Bölüm~\ref{sec:tartisma}
yorumları ve sınırlılıkları sunar; Bölüm~\ref{sec:sonuc} çalışmayı bağlar.
